\section{Liniowe operatory w przestrzeni euklidesowej}

Rozpatrzmy $n$-wymiarową przestrzeń euklidesową $V$. Niech będzie dany dowolny operator liniowy $A$ w tej przestrzeni.

\begin{definition}
    Operator liniowy $A^T$ nazywa się \textbf{operatorem transponowanym do operatora} $A$, jeśli dla dowolnych dwóch wektorów $x, y \in V$:
    \begin{equation}\label{eq:def_transposed_operator}
        \inner{Ax}{y} = \inner{x}{A^Ty}
    \end{equation}
\end{definition}
\begin{remark}
Odpowiednik operatora transponowanego w przestrzeni unitarnej jest operator sprzężony.
\end{remark}
Następujące twierdzenia podamy bez dowodów, ponieważ są analogiczne jak w przypadku dla przestrzeni unitarnej.

\begin{theorem}
Dla każdego liniowego operatora $A$ w przestrzeni euklidesowej $V$ zawsze istnieje operator transponowany $A^T$ i to tylko jeden.
\end{theorem}
% \begin{proof}
% Patrz Twierdzenie \ref{theorem:existing_of_conjugated_operator}
% \end{proof}

\begin{theorem}
    W przestrzeni euklidesowej $V$ z bazą ortonormalną, transponowanemu operatorowi $A^T$ odpowiada macierz transponowana.
\end{theorem}

\begin{theorem}\label{theorem:properties_of_transposed_operator}
    Z definicji operatora transponowanego mamy następujące własności:
    \begin{enumerate}[leftmargin=1cm, itemsep=2pt, label=(\arabic*)]
        \item $(A^T)^T = A$
        \item $(A + B)^T = A^T + B^T$
        \item $(\lambda A)^T = \lambda(A^T)$, gdzie $\lambda \in \R$
        \item $(AB)^T = B^TA^T$
    \end{enumerate}
\end{theorem}

\begin{definition}\label{def:real_normal_operator}
Operator liniowy $A$ nazywa się \textit{normalnym}, jeśli:
$$AA^T = A^TA$$
\end{definition}
\begin{remark}
W przestrzeni unitarnej, operator normalny definiujemy jako:
$$A^*A = AA^*,$$ co pokrywa się z $$AA^T = A^TA$$ w przestrzeni euklidesowej.
\end{remark}
\begin{theorem}
    Operator normalny $A$ w przestrzeni euklidesowej $V$ z bazą ortonormalną ma macierz normalną.
\end{theorem}

\begin{definition}\label{def:symmetric_operator}
Operator liniowy $S$ nazywa się \textit{symetrycznym}, jeśli:
$$S = S^T$$
\end{definition}
\begin{theorem}\label{theorem:sym_op_has_same_matrix}
    Operator symetryczny $S$ w przestrzeni euklidesowej $V$ z bazą ortonormalną ma macierz symetryczną.
\end{theorem}

\begin{remark}
Odpowiednik operatora symetrycznego w przestrzeni unitarnej jest operator hermitowski.
\end{remark}

\begin{definition}\label{def:skew_symmetric_operator}
Operator liniowy $K$ nazywa się \textit{antysymetrycznym}, jeśli:
$$K = -K^T$$
\end{definition}

\begin{theorem}
    Operator antysymetryczny $K$ w przestrzeni euklidesowej $V$ z bazą ortonormalną ma macierz antysymetryczną.
\end{theorem}
\begin{proof}
    Z definicji operatora antysymetrycznego dla dowolnych dwóch wektorów $\textbf{e}_i$, $\textbf{e}_k$ z bazy ortonormalnej $\base{e}$, gdzie $i, k \in \{1, \dots, n\}$:
    $$\inner{K\textbf{e}_k}{\textbf{e}_i} = \inner{\textbf{e}_k}{-K\textbf{e}_i}$$
    Korzystając ze wzoru \eqref{eq:def_ortonormal_coeffs} oraz z jednorodności i przemienności iloczynu skalarnego:
    $$ a_{ik} = \inner{K\textbf{e}_k}{\textbf{e}_i} = \inner{\textbf{e}_k}{-K\textbf{e}_i} = \inner{-K\textbf{e}_i}{\textbf{e}_k} = -\inner{K\textbf{e}_i}{\textbf{e}_k} = -a_{ki}, $$ gdzie $i, k \in \{1, \dots, n\}$. Z tego wynika, że operatorowi antysymetrycznemu odpowiada macierz antysymetryczna.
\end{proof}

\begin{definition}\label{def:ortogonal_operator}
    Operator liniowy $O$ nazywa się \textit{ortogonalnym} lub inaczej \textit{izometrią liniową}, jeśli zachowuje metrykę przestrzeni, tj. jeśli dla dowolnych dwóch wektorów $x, y \in V$:
    \begin{equation}\label{eq:def_ortogonal_operator}
        \inner{Ox}{Oy} = \inner{x}{y}
    \end{equation}
\end{definition}
Równość \eqref{eq:def_ortogonal_operator}, jak i dla operatora unitarnego, możemy przepisać w następujący sposób na mocy \eqref{eq:def_transposed_operator}:
$$\inner{Oa}{b} = \inner{a}{O^Tb}$$
Podstawiając pod $a = x$, a pod $b = Oy$ otrzymujemy:
$$\inner{Ox}{Oy} = \inner{x}{O^TOy} = \inner{x}{y}$$
Z tego wynika:
\begin{equation}\label{eq:def_ortogonal_operator_2}
    O^TO = I,
\end{equation} gdzie $I$ operator tożsamościowy (macierz jednostkowa). 
\begin{remark}
Odpowiednik operatora ortogonalnego w przestrzeni unitarnej jest operator unitarny.
\end{remark}
\begin{theorem}
    Operator ortogonalny $O$ w przestrzeni euklidesowej $V$ z bazą ortonormalną ma macierz ortogonalną.
\end{theorem}

Operatory symetryczne, antysymetryczne oraz ortogonalne, jak i dla przestrzeni unitarnej, są szczególnymi przypadkami operatora normalnego.

Jak widzimy z powyższego, utożsamiamy macierze hermitowskie, unitarne, symetryczne, antysymetryczne oraz ortogonalne z operatorami liniowymi lub inaczej, w przestrzeni unitarnej każda z tych macierzy reprezentuje przekształcenie liniowe.
